\section*{Bab \ref{ch:b3:complex-kinetics} --- Kinetika Kompleks: Rantai, Enzim, dan Osilasi}

\begin{solution}{exo:b3:complex-kinetics:1}
$\ce{Cl2 -> 2 Cl}$: inisiasi; $\ce{Cl + H2 -> HCl + H}$ dan $\ce{H + Cl2 -> HCl + Cl}$:
propagasi; $\ce{2 Cl -> Cl2}$: terminasi. Pembawa: Cl dan H. Reaksi keseluruhan
(jumlah kedua tahap propagasi): $\ce{H2 + Cl2 -> 2 HCl}$.
\end{solution}

\begin{solution}{exo:b3:complex-kinetics:2}
$[\mathrm M]_{1/2} = k_2/k_{-1} = \qty{1.0e-3}{mol/L} = \qty{1.0}{mol/m^3}$; $p = cRT = 1.0 \times 8.314 \times 750 =
\qty{6.2}{kPa}$.
\end{solution}

\begin{solution}{exo:b3:complex-kinetics:3}
$v/V_{\max} = [\mathrm S]/(K_M + [\mathrm S])$: $\frac12$, $\frac34$, $\frac{9}{10}$. 90\,\% dari $V_{\max}$ memerlukan $[\mathrm S] = 9K_M$.
\end{solution}

\begin{solution}{exo:b3:complex-kinetics:4}
$k_{\mathrm{cat}}/K_M = 100/\num{0.80e-3} = \qty{1.25e5}{L.mol^{-1}.s^{-1}}$: sekitar 60\,000 kali
di bawah batas difusi, dan dekat dengan nilai khasnya.
\end{solution}

\begin{solution}{exo:b3:complex-kinetics:5}
Pendekatan keadaan tunak untuk \ce{CH3CO} memberikan $k_2[\ce{CH3}][\ce{CH3CHO}] = k_3[\ce{CH3CO}]$. Keadaan
tunak untuk \ce{CH3}:
$k_1[\ce{CH3CHO}] - k_2[\ce{CH3}][\ce{CH3CHO}] + k_3[\ce{CH3CO}] - 2k_4[\ce{CH3}]^2 = 0$.
Dengan menjumlahkan: $k_1[\ce{CH3CHO}] =
2k_4[\ce{CH3}]^2$, sehingga $[\ce{CH3}] = (k_1/2k_4)^{1/2}[\ce{CH3CHO}]^{1/2}$, dan $\dd[\ce{CH4}]/\dd t = k_2[\ce{CH3}][\ce{CH3CHO}] =
k_2(k_1/2k_4)^{1/2}[\ce{CH3CHO}]^{3/2}$.
\end{solution}

\begin{solution}{exo:b3:complex-kinetics:6}
Pada setengah konversi $[\ce{H2}] = [\ce{Br2}] = \frac12$, $[\ce{HBr}] = 1$ (satuan
relatif): $v/v_0 = \frac12 \times (\frac12)^{1/2}/(1 + 0.10 \times 1/0.5) = 0.354/1.2 = 0.29$.
Tanpa inhibisi nilainya 0.354: inhibisi membagi lajunya lebih lanjut dengan 1.2.
\end{solution}

\begin{solution}{exo:b3:complex-kinetics:7}
Ledakan terjadi jika $1000p > 6000/p + 10p^2$, yaitu $10p^3 - 1000p^2 + 6000 < 0$, yang
akar-akar positifnya 2.48 dan \qty{99.9}{kPa}: campuran itu meledak antara
sekitar 2.5 dan \qty{100}{kPa} (batas pertama dan kedua model itu).
\end{solution}

\begin{solution}{exo:b3:complex-kinetics:8}
$1/v = 1/V_{\max} + (K_M/V_{\max})/[\mathrm S]$: $5.0 = 1/V_{\max} + 2.0K_M/V_{\max}$ dan $2.5 = 1/V_{\max} + 0.5K_M/V_{\max}$.
Dengan mengurangkan: $K_M/V_{\max} = \qty{1.667}{s.mM.\micro M^{-1}}$, lalu $1/V_{\max} = 1.667$:
$V_{\max} = \qty{0.60}{\micro M.s^{-1}}$, $K_M = \qty{1.0}{mM}$.
\end{solution}

\begin{solution}{exo:b3:complex-kinetics:9}
Perpotongan yang sama pada sumbu $1/v$: kompetitif. $K_M^{\mathrm{app}} = 1/0.417 = \qty{2.40}{mM} = 3K_M$,
sehingga $1 + 50/K_i = 3$ dan $K_i = \qty{25}{\micro M}$.
\end{solution}

\begin{solution}{exo:b3:complex-kinetics:10}
$N = \qty{1.01}{mol/L}$: $t_{1/2} = \ln(1.01/0.010 - 1)/(0.50 \times 1.01) = \ln100/0.505 = \qty{9.1}{s}$.
Laju $kx(N - x)$ paling besar pada $x = N/2$: saat yang sama, titik belok
kurva S itu.
\end{solution}

\begin{solution}{exo:b3:complex-kinetics:11}
Trace $B - 2$, determinan 1. $B = 1.5$: $\lambda = -0.25 \pm 0.968\iu$, fokus stabil
(osilasi teredam menuju keadaan tunak). $B = 2.5$: $\lambda = 0.25 \pm 0.968\iu$,
fokus tidak stabil: osilasinya membesar sampai mencapai \omterm{def:b3:complex-kinetics:oscillating}{siklus limit}.
\end{solution}

\begin{solution}{exo:b3:complex-kinetics:12}
$1/\tau = \num{1.0e8} \times \num{5.40e-5} + \num{1.0e3} = \qty{6.4e3}{s^{-1}}$, $\tau = \qty{156}{\micro s}$.
Ukurlah $\tau$ pada beberapa konsentrasi total: $1/\tau$ terhadap
$[\mathrm A]_e + [\mathrm B]_e$ adalah garis lurus dengan kemiringan $k_1$ dan
perpotongan $k_{-1}$.
\end{solution}

\begin{solution}{pb:b3:complex-kinetics:1}
\textbf{1.} Pada awal reaksi, $[\mathrm S]$ diketahui dan produknya, yang dapat
menginhibisi atau bereaksi balik, belum ada.
\textbf{2.} Titik-titik dengan $[\mathrm S]$ rendah, pada $1/[\mathrm S]$ dan $1/v$ yang
besar: titik-titik itu menentukan kemiringan, dan pembalikan memperbesar
galatnya.
\textbf{3.} Pencocokan itu memberi bobot pada setiap laju terukur sebagaimana
diukur; garis Lineweaver--Burk memberikan parameter yang terbias.
\textbf{4.} $K_M$: $0.773$ terletak 1.2 galat baku di bawah 0.801; $V_{\max}$: $0.491$
terletak 2.2 galat baku di bawah 0.502. Keduanya rendah, sesuai dengan bias
yang diharapkan.
\textbf{5.} $k_{\mathrm{cat}} = 0.502/0.0050 = \qty{100}{s^{-1}}$.
\textbf{6.} $100/\num{0.801e-3} = \qty{1.25e5}{L.mol^{-1}.s^{-1}}$.
\textbf{7.} Sekitar 60\,000 kali di bawah batas difusi; \omterm{def:b3:complex-kinetics:enzyme}{enzim} yang khas.
\textbf{8.} Kompetitif: $V_{\max}$ tidak berubah, $K_M$ semu tumbuh bersama $[\mathrm I]$.
\textbf{9.} $K_M^{\mathrm{app}} = K_M(1 + [\mathrm I]/K_i) = K_M + (K_M/K_i)[\mathrm I]$.
\textbf{10.} $K_i = 0.799/0.0321 = \qty{24.9}{\micro M}$.
\textbf{11.} $24.9 \pm 4.30 \times 0.9 = 24.9 \pm \qty{3.9}{\micro M}$.
\textbf{12.} Tetapan disosiasi kompleks enzim--inhibitor: konsentrasi inhibitor
yang menempati setengah \omterm{def:b3:complex-kinetics:enzyme}{enzim} bebas.
\textbf{13.} $v_i/v_0 = (K_M + [\mathrm S])/(K_M(1 + [\mathrm I]/K_i) + [\mathrm S])$.
\textbf{14.} $2/3 = 0.67$.
\textbf{15.} $2/12 = 0.17$.
\textbf{16.} $11/21 = 0.52$: pada konsentrasi substrat yang tinggi, substrat
mengalahkan inhibitor dalam persaingan.
\textbf{17.} Pada $[\mathrm S] = K_M$, $v_i/v_0 = 2/(2 + [\mathrm I]/K_i) = 0.10$ memberikan
$[\mathrm I]/K_i = 18$.
\textbf{18.} $18 \times 24.9 = \qty{448}{\micro M}$.
\textbf{19.} Tahap pertama cepat: setiap molekul \omterm{def:b3:complex-kinetics:enzyme}{enzim} dengan cepat melepaskan
satu $\mathrm P_1$ dan terperangkap sebagai \omterm{def:b3:complex-kinetics:enzyme}{asil-enzim}; sesudahnya $\mathrm P_1$
hanya muncul secepat tahap kedua yang lambat membebaskan \omterm{def:b3:complex-kinetics:enzyme}{enzim}.
\textbf{20.} $1.28/2.0 = 0.64 = (k_2/(k_2 + k_3))^2$, sehingga $k_2/(k_2 + k_3) = 0.80$.
\textbf{21.} $200/2.0 = \qty{100}{s^{-1}}$, sama dengan pada Bagian I.
\textbf{22.} $k_2 = 0.80 \times 625 = \qty{500}{s^{-1}}$, $k_3 = \qty{125}{s^{-1}}$.
\textbf{23.} $500 \times 125/625 = \qty{100}{s^{-1}}$.
\textbf{24.} \textbf{$[\mathrm I]_{90\,\%} = 18K_i \approx \qty{0.45}{mM}$ pada $[\mathrm S] = K_M$.}
\end{solution}
